% Source: Solutions/Ansh-PS1-Q1-5.pdf, p. 5 (typed). Restated problem statement omitted.
\submittedby{Ansh}

\begin{proof}
Let $(X,\tau)$ be a topological space and let $A\subseteq X$. We prove both directions of the
biconditional statement.

\textbf{Forward Direction ($\Longrightarrow$).} Suppose that $A$ is open in $X$, so $A\in\tau$.
Let $x\in A$ be arbitrary. Choose $U:=A$. Since $A$ is open in $X$, $U$ is an open set in $X$,
and by reflexivity of set inclusion ($A\subseteq A$), we immediately have
\[ x\in U\subseteq A. \]
Thus, for each $x\in A$, there exists an open set $U$ in $X$ such that $x\in U\subseteq A$.

\textbf{Reverse Direction ($\Longleftarrow$).} Conversely, suppose that for every $x\in A$,
there exists an open set $U_x\in\tau$ such that $x\in U_x\subseteq A$. We claim that
\[ A=\bigcup_{x\in A}U_x. \]
We establish this set equality by proving mutual inclusion:
\begin{enumerate}
  \item \emph{Proof that $A\subseteq\bigcup_{x\in A}U_x$:} Let $y\in A$ be arbitrary. By
        hypothesis, the corresponding open set $U_y$ satisfies $y\in U_y$. Since $U_y$ is one
        of the sets in the indexed family $\{U_x\}_{x\in A}$, it follows from the definition
        of the union that $y\in\bigcup_{x\in A}U_x$. Hence, $A\subseteq\bigcup_{x\in A}U_x$.
        (Note that if $A=\emptyset$, the empty union is $\emptyset$, so the inclusion holds
        vacuously.)
  \item \emph{Proof that $\bigcup_{x\in A}U_x\subseteq A$:} Let $y\in\bigcup_{x\in A}U_x$ be
        arbitrary. By definition of the union, there exists some index $x_0\in A$ such that
        $y\in U_{x_0}$. By hypothesis, $U_{x_0}\subseteq A$, which implies $y\in A$. Hence,
        $\bigcup_{x\in A}U_x\subseteq A$.
\end{enumerate}
Combining both inclusions yields $A=\bigcup_{x\in A}U_x$. Because each $U_x$ is an open set in
$X$ and a topology $\tau$ is closed under arbitrary unions, we conclude that
$A=\bigcup_{x\in A}U_x\in\tau$; that is, $A$ is open in $X$.
\end{proof}
